\documentclass[10pt,a4paper]{amsart}
\usepackage[margin=19mm]{geometry}
\usepackage{amsmath,amssymb,mathtools}
\usepackage[scr=rsfs,cal=euler]{mathalfa}
\usepackage{xfrac}
\usepackage[unicode,hidelinks,bookmarks=false]{hyperref}
\newcommand{\ie}{i.e.\ }
\newcommand{\cf}{cf.\ }
\newcommand{\resp}{respectively\ }
\newcommand{\Subsection}{\subsection}
\providecommand{\nobreakdash}{\nobreak\hbox{-}}
\setlength{\emergencystretch}{2em}
\multlinegap0pt
\usepackage{bm}
\usepackage{bbm}
\usepackage{dsfont}
\usepackage{xspace}
\usepackage{cancel}
\usepackage{mathtools}
\usepackage{cleveref}
\usepackage{amsthm}
\makeatletter
\@ifundefined{theorem}{\newtheorem{theorem}{Theorem}}{}
\@ifundefined{lemma}{\newtheorem{lemma}[theorem]{Lemma}}{}
\makeatother
\def\acc#1{\left\{ #1 \right\}}
\renewcommand{\le}{\leqslant}
\title[More characterizations of morphic words]{More characterizations of morphic words}
\author{Golnaz Badkobeh, Pascal Ochem}
\date{Source: arXiv:2312.10757v2 (CC BY 4.0)}
\begin{document}
\maketitle

\noindent\emph{Source and licence.} More characterizations of morphic words. Version arXiv:2312.10757v2 (CC BY 4.0), distributed under the Creative Commons Attribution 4.0 International licence, as recorded in the arXiv licence metadata for this exact version and shown on the abstract page https://arxiv.org/abs/2312.10757v2. Full rights notice: licenses/arxiv-2312.10757-CC-BY-4.0.txt.\par\smallskip

\noindent\textit{Editorial scope.} Counted unit: the Theorem whose source label is \texttt{None} in the article (source0.tex lines 267--270, source label \texttt{None}), delivered together with the article's own complete original proof of it (source0.tex lines 272--297, one \texttt{proof} environment of 26 lines). No mathematics is written, repaired or re-derived: statement and proof are the article's own text line for line. Required context retained uncounted: lem:b3 (lines 184--193). Every definition, display and label of those lines is retained unchanged. Omitted from this package and not redistributed: 1--183, 194--266, 271--271, 298--885. Nothing inside a retained excerpt is omitted or altered: every retained body is delivered exactly as the article's lines are written, so no omission receipt of any kind accompanies this package. Citations: the retained excerpts carry no citation command at all, so no bibliography entry of the article is transcribed and no reference is redistributed. Reference adaptation: none was needed and none was made; every reference inside the delivered unit resolves inside it, so no unresolved reference or citation is produced. Printed slips kept literal: no printed word, symbol, hypothesis or number of the counted statement or of its proof was altered. Layout and class: the article's own class is \texttt{article} with packages including inputenc, amsmath, amssymb, amsthm, graphicx, hyperref, url, cleveref, subcaption, todonotes, enumerate; the wrapper is the lane's pinned \texttt{amsart} editorial wrapper at 10pt on A4, which restates the theorem-like environments the retained text needs and adds the article's own macro definitions that the retained text needs (\texttt{\textbackslash acc}, \texttt{\textbackslash le}), with extra wrapper packages (bm, bbm, dsfont, xspace, cancel, mathtools, cleveref). The delivered numbering is the unit's own. Assets: no retained span contains a figure environment, a graphics-inclusion command, a PDF-page inclusion, a table or a TikZ picture; no image file is copied into this package, the package declares no asset dependency and no image file is redistributed with it. Licence: version arXiv:2312.10757v2 of this article is distributed under the Creative Commons Attribution 4.0 International licence, as recorded in the arXiv licence metadata for this exact version and shown on the abstract page https://arxiv.org/abs/2312.10757v2. A full rights notice is delivered at licenses/arxiv-2312.10757-CC-BY-4.0.txt. Characters: every retained excerpt is pure ASCII, so the delivered engine, XeLaTeX with its default Latin Modern text font, reports no missing character.
\begin{lemma}\label{lem:b3}
Let $w$ be a bi-infinite word and let $m:\Sigma^*_3\rightarrow\Sigma^*_k$ be such that 
\begin{itemize}
 %\item $w$ has the same factor set of length $5\times\max\acc{|m(\textnormal{\texttt{0}})|,|m(\textnormal{\texttt{1}})|,|m(\textnormal{\texttt{2}})|}$ as $m(b_3)$,
 \item $w$ is in $\acc{m(\textnormal{\texttt{0}}),m(\textnormal{\texttt{1}}),m(\textnormal{\texttt{2}})}^\omega$,
 \item $m(\textnormal{\texttt{0}})=axb$ and $m(\textnormal{\texttt{1}})=ab$,
 \item $w$ avoids $SQ_t$ with $t\le\min\acc{|m(\textnormal{\texttt{0}})|,|m(\textnormal{\texttt{1}})|,|m(\textnormal{\texttt{2}})|}$.
\end{itemize}
then $w$ has the same factor set as $m(b_3)$.
\end{lemma}

\begin{theorem}
Every binary bi-infinite word containing only the squares \textnormal{\texttt{00}}, \textnormal{\texttt{11}}, \textnormal{\texttt{0000}}, \textnormal{\texttt{0001100011}}, and \textnormal{\texttt{1000010000}}
has the same factor set as $g_5(b_3)$.
\end{theorem}

\begin{proof}
We say that a binary word, finite or infinite, is \emph{good} if it contains
no squares other than these five squares.
It is easy to check that a bi-infinite word is good if and only if it avoids $SQ_6$ and\\
$\{\texttt{101},\texttt{1001},\texttt{1111},\texttt{000000},\texttt{010001},
\texttt{100010},\texttt{0000010},\texttt{0100001},\texttt{1000110},\texttt{1110001}\}$.

First, we show that $g_5(b_3)$ is good.
A computer check shows that the factors of length 200 of $g_5(b_3)$ are good.
Now we have to show that $g_5(b_3)$ also avoids larger squares.
Unfortunately, checking directly that $g_5(b_3)$ avoids $SQ_6$ with Rosenfeld's program
is out of reach because 6 variables is too many for a morphism as large as $g_5$.
Thus, we had to resort to the following workaround.
Every factor of length 29 of $g_5(b_3)$ contains the factor \texttt{00000}.
So if $g_5(b_3)$ contains a square with period at least 31, then $g_5(b_3)$ contains an occurrence
of the pattern $P=ABBBBBCABBBBBC$ (such that $B\mapsto\texttt{0}$).
However, we have checked that $g_5(b_3)$ avoids~$P$. So $g_5(b_3)$ avoids $SQ_6$.

We construct the set $S_5^{100}$ defined as follows:
a word $v$ is in $S_5^{100}$ if and only if there exists a good word $pvs$ such that $|p|=|v|=|s|=100$.
We check that $S_5^{100}$ is exactly the set of factors of length 100 of $g_5(b_3)$.
%This implies that every bi-infinite good word $r_2$ has the same set of factors of length 100 as $g_5(b_3)$.
Using the same argument as in the previous proof, we obtain that every bi-infinite good word $r_2$
is in $\acc{g_5(\texttt{0}),g_5(\texttt{1}),g_5(\texttt{2})}^\omega$.
Then $r_2$ has the same factor set as $g_5(b_3)$ by~\Cref{lem:b3}.
\end{proof}

\end{document}
